\chapter{Conclusiones y trabajos futuros}
\label{ch:conclusiones}

%\avisoLocalizacionArchivo 

\begin{comment}
Las conclusiones deben cerrar el documento, destacando los aspectos más importantes de la ejecución del \gls{TFG}.  Debe analizar qué objetivos se han alcanzado y en qué grado, qué objetivos se han tenido que dejar fuera del proyecto y por qué, y qué líneas de trabajo futuro abre el \gls{TFG}. 

Fíjate en que los objetivos abren el trabajo personal y las conclusiones lo cierran.  Procura mantener un orden que resalte esta relación, pero no te limites a parafrasear los objetivos.
\end{comment}

Este trabajo partió de una doble motivación: resolver maniobras orbitales que no admiten una fórmula cerrada y acercar su planificación a quien no es especialista. El resultado es un sistema completo que recorre todo el camino, desde un entorno físico contrastado con datos de misiones reales hasta una interacción en lenguaje natural, pasando por un conjunto de agentes que aprenden las maniobras y por los solucionadores clásicos allí donde el óptimo ya se conoce. Las páginas anteriores han descrito cada pieza; este capítulo valora en qué medida se alcanzaron las metas fijadas en el Capítulo~\ref{ch:objetivos}, qué quedó deliberadamente fuera y qué caminos abre el trabajo.

\section{Conclusiones}

La base física del proyecto no solo se construyó, sino que su exigencia de validación resultó más fértil de lo previsto. Al contrastar el modelo atmosférico con datos in-situ de las sondas se detectaron y corrigieron errores que, de otro modo, habrían contaminado todo lo demás: el más llamativo, una densidad de referencia de Neptuno equivocada en un factor 100. Esa misma disciplina llevó a desconfiar de las cifras no trazables ---una asistente generalista llegó a errar en un factor 2500 un dato de Júpiter--- y a apoyarse únicamente en fuentes primarias. El entorno resultante reproduce los datos de referencia con un factor de desviación inferior a 2, coincide con la teoría de Brouwer con un error menor del 1\,\% en los cuerpos de achatamiento moderado y supera la totalidad de una batería de 54 pruebas automáticas. La validación se convirtió de esta forma en una fuente de aportaciones.

Sobre esa base, los agentes de \gls{RL} cumplieron sus objetivos con holgura y, en varios casos, los superaron. El primero demostró que la tubería de aprendizaje funciona, aprendiendo la transferencia de Hohmann sin conocer la fórmula hasta quedar a un 0,03\,\% del óptimo; pero el hallazgo relevante fue el siguiente, cuando la adimensionalización del problema permitió que un único agente resolviera cualquier transferencia en cualquier cuerpo sin reentrenar, con un exceso inferior al 1\,\%, gracias a la invariancia de escala de la maniobra. La extensión tridimensional incorporó el cambio de plano manteniendo el exceso por debajo del 2\,\%. El aerofrenado abordó, ya sí, una maniobra sin solución analítica: un agente por cuerpo aprendió a apurar la atmósfera hasta el límite seguro de presión dinámica, resolviendo 50 de 50 escenarios en cada uno de los siete planetas con atmósfera. Y el mantenimiento orbital se completó dentro del plazo: mantiene la órbita en banda en 30 de 30 escenarios y en los siete cuerpos, confirmando además, con honestidad, que el coste de combustible de esta maniobra lo fija esencialmente la física.

La capa de orquestación culminó el sistema. Un \gls{LLM} interpreta las peticiones en lenguaje natural, invoca el agente o el solucionador que corresponde ---mediante el patrón de \textit{tool use}--- y explica los resultados sin inventar ninguna cifra, pues se le ancla a los datos que devuelven las herramientas. Se resolvió, siguiendo la recomendación de la dirección, con un modelo abierto y gratuito, lo que cumple el alcance de demostrador previsto y lo hace sin coste. Por último, la evaluación se ejerció de forma transversal a lo largo del trabajo: cada maniobra se comparó con su referencia clásica y sus limitaciones se documentaron sin adornos.

%\section{Aportaciones principales}

Más allá del cumplimiento de cada objetivo, el trabajo deja algunas contribuciones que conviene destacar. La primera es metodológica: una validación cuantitativa frente a datos de misión que no se limita a confirmar, sino que descubre y corrige errores, incluida la reconstrucción de las atmósferas de Urano y Neptuno a falta de mediciones directas. La segunda es la generalización de las maniobras mediante análisis dimensional, que convierte un agente entrenado una sola vez en una herramienta válida para todo el Sistema Solar. La tercera es la resolución, con aprendizaje, de maniobras que ninguna fórmula captura ---el aerofrenado y el mantenimiento---, apoyadas en un modelo físico realista. Y, de manera transversal, una honestidad metodológica sostenida: se documentan los supuestos del modelo, se admite dónde una maniobra es viable solo en simulación y se distingue con claridad lo demostrado de lo estimado.

%\section{Limitaciones y decisiones de alcance}

Algunas metas se acotaron a conciencia. La capa \gls{LLM} se planteó desde el principio como un demostrador de la arquitectura y no como un producto; su uso intensivo está limitado por la cuota gratuita del servicio empleado, suficiente para probar el concepto pero no para una explotación a gran escala. Las transferencias interplanetarias se resolvieron con métodos clásicos y no con aprendizaje, por la razón deliberada de no sustituir con \gls{RL} aquello que ya tiene un óptimo conocido y fiable. En el mantenimiento orbital se corrige la geometría de la órbita frente al arrastre, pero no se combate la lenta rotación del plano que induce el achatamiento: hacerlo exigiría maniobras fuera del plano de coste prohibitivo, que en la práctica no se realizan. Y tanto en el aerofrenado como en el mantenimiento se optó por un especialista por cuerpo en lugar de un único agente universal, porque cada atmósfera tiene su propia escala. Ninguna de estas decisiones responde a un fracaso, sino a un criterio de realismo y de honestidad con el alcance del trabajo.

\section{Trabajos futuros}

El trabajo abre varias continuaciones naturales. La más inmediata es un mantenimiento orbital en tres dimensiones que gestione también la orientación del plano y la inclinación, con control fuera del plano, reutilizando el propagador ya validado. Una segunda línea es sustituir los especialistas por cuerpo por un agente único multicuerpo, incorporando al estado los parámetros que distinguen cada atmósfera. Cabría asimismo extender los agentes a las transferencias de bajo empuje, de control continuo, y enriquecer el diseño interplanetario con asistencias gravitatorias, que abaratan los saltos entre planetas muy separados. Del lado de la orquestación, desplegar la capa \gls{LLM} sobre un modelo ejecutado localmente eliminaría la dependencia de una cuota externa y permitiría ampliar el catálogo de herramientas. Y, en el
plano físico, un modelo atmosférico dependiente de la latitud y de la actividad solar elevaría la fidelidad de las maniobras más sensibles al arrastre.

%\section{Reflexión final}

Los objetivos abrieron este trabajo comprometiéndose con una idea doble: emplear el aprendizaje donde la física carece de fórmula y hacer accesible la planificación de maniobras. El recorrido lo confirma. Los agentes resuelven, aprendiendo, lo que no se sabe resolver de otro modo ---el aerofrenado, el mantenimiento---, se apoyan en una física que se ha contrastado y no supuesto, y quedan al alcance de cualquiera a través de una conversación en lenguaje natural. Que un mismo sistema abarque desde la densidad de la atmósfera de Neptuno hasta la frase con la que un usuario pide llegar a una órbita es, quizá, la mejor síntesis de lo aprendido: que la utilidad de una herramienta nace tanto del rigor con que se construye como de la sencillez con que se ofrece.